\chapter{Hito 5: Construcción de la Capa Staging (Bronze)}

\section{Objetivo}

El objetivo de este hito es construir la primera capa física del Data Warehouse.

La capa staging será el punto de entrada de todos los datos al sistema.

Al finalizar este hito el estudiante será capaz de:

\begin{itemize}
\item Crear schemas en PostgreSQL.
\item Diseñar tablas de aterrizaje.
\item Comprender la función de la capa Bronze.
\item Implementar una estrategia de carga desacoplada de las transformaciones.
\item Preparar el sistema para recibir datos desde Python.
\end{itemize}

\section{Resultado esperado}

Al finalizar este hito deberá existir la siguiente estructura lógica:

\begin{verbatim}
PostgreSQL

+-- staging
|
+-- core
|
+-- mart
\end{verbatim}

Y dentro del schema staging:

\begin{verbatim}
staging.tipo_cambio_raw

staging.balance_bancario_raw
\end{verbatim}

\section{¿Qué es la capa Staging?}

La capa staging es el área donde aterrizan los datos provenientes de las fuentes externas.

Su objetivo es almacenar los datos en su forma más cercana al origen.

No se realizan:

\begin{itemize}
\item Cálculos.
\item Transformaciones complejas.
\item Generación de indicadores.
\item Limpieza profunda.
\end{itemize}

La regla general es:

\begin{quote}
Cargar primero. Transformar después.
\end{quote}

\section{Arquitectura del Data Warehouse}

La arquitectura que construiremos será:

\begin{verbatim}
Fuentes
|
V
Staging (Bronze)
|
V
Core (Silver)
|
V
Mart (Gold)
\end{verbatim}

\section{Paso 1: Crear carpeta para scripts SQL}

Verificar:

\begin{verbatim}
sql/
\end{verbatim}

Crear:

\begin{verbatim}
sql/01_create_schemas.sql
\end{verbatim}

\section{Paso 2: Crear los schemas}

Agregar:

\begin{lstlisting}[language=SQL]
CREATE SCHEMA IF NOT EXISTS staging;

CREATE SCHEMA IF NOT EXISTS core;

CREATE SCHEMA IF NOT EXISTS mart;
\end{lstlisting}

Guardar.

\section{Paso 3: Ejecutar el script}

Desde psql:

\begin{lstlisting}[language=SQL]
\i sql/01_create_schemas.sql
\end{lstlisting}

Verificar:

\begin{lstlisting}[language=SQL]
SELECT schema_name
FROM information_schema.schemata
WHERE schema_name IN
(
'staging',
'core',
'mart'
);
\end{lstlisting}

Resultado esperado:

\begin{verbatim}
staging
core
mart
\end{verbatim}

\section{Paso 4: Crear script de tablas staging}

Crear:

\begin{verbatim}
sql/02\_staging\_ddl.sql
\end{verbatim}

\section{Paso 5: Diseñar tabla tipo\_cambio\accentedoperators                                                                                                                                                                                                                                                                                                                                                                                                                                                                                                                                                                                                                                                                                                                                                                                                                                                                                                                                                                                                                                                                                                                                                                                                                                                                                                                                                                                                                                                                                                                                                                                                                                                                                                                                                                                                                                                                                                                                                                                                                                                                                                                                                                                                                                                                                                                                                                                                                                                                                                                                                                                                                                                                                                                                                                                                                                                                                                                                                                                                                                                                                                                                                                             _raw}

Agregar:

\begin{lstlisting}[language=SQL]
CREATE TABLE IF NOT EXISTS
staging.tipo_cambio_raw
(
fecha TEXT,

```
moneda TEXT,

tasa_compra TEXT,

tasa_venta TEXT,

archivo_origen TEXT,

cargado_en TIMESTAMP
    DEFAULT CURRENT_TIMESTAMP
```

);
\end{lstlisting}

\section{¿Por qué usamos TEXT?}

Esta es una práctica común en Data Engineering.

Ventajas:

\begin{itemize}
\item La carga nunca falla por formato.
\item Se preserva el dato original.
\item La validación ocurre posteriormente.
\item Se facilita el diagnóstico de errores.
\end{itemize}

\section{Paso 6: Diseñar tabla balance_bancario_raw}

Agregar:

\begin{lstlisting}[language=SQL]
CREATE TABLE IF NOT EXISTS
staging.balance_bancario_raw
(
periodo TEXT,

```
entidad_cod TEXT,

entidad_nombre TEXT,

tipo_entidad TEXT,

activos_totales TEXT,

cartera_creditos TEXT,

depositos TEXT,

patrimonio TEXT,

resultado_neto TEXT,

depositos_me TEXT,

cartera_me TEXT,

archivo_origen TEXT,

cargado_en TIMESTAMP
    DEFAULT CURRENT_TIMESTAMP
```

);
\end{lstlisting}

\section{Paso 7: Ejecutar el DDL}

Ejecutar:

\begin{lstlisting}[language=SQL]
\i sql/02_staging_ddl.sql
\end{lstlisting}

\section{Paso 8: Verificar tablas creadas}

Consultar:

\begin{lstlisting}[language=SQL]
SELECT table_name
FROM information_schema.tables
WHERE table_schema='staging';
\end{lstlisting}

Resultado esperado:

\begin{verbatim}
tipo_cambio_raw

balance_bancario_raw
\end{verbatim}

\section{Paso 9: Inspeccionar estructura}

Consultar:

\begin{lstlisting}[language=SQL]
SELECT
column_name,
data_type
FROM information_schema.columns
WHERE table_schema='staging'
AND table_name='tipo_cambio_raw';
\end{lstlisting}

Verificar que todas las columnas fueron creadas correctamente.

\section{Paso 10: Crear script de prueba}

Crear:

\begin{verbatim}
sql/test_staging.sql
\end{verbatim}

Contenido:

\begin{lstlisting}[language=SQL]
SELECT *
FROM staging.tipo_cambio_raw;

SELECT *
FROM staging.balance_bancario_raw;
\end{lstlisting}

Actualmente ambas tablas deben estar vacías.

\section{Paso 11: Documentar la capa staging}

Crear:

\begin{verbatim}
docs/staging_design.md
\end{verbatim}

Documentar:

\begin{itemize}
\item Nombre de tabla.
\item Fuente original.
\item Frecuencia.
\item Número esperado de columnas.
\item Responsable de carga.
\end{itemize}

Ejemplo:

\begin{verbatim}
Tabla:
staging.tipo_cambio_raw

Fuente:
Banco Central

Frecuencia:
Diaria

Objetivo:
Almacenar datos crudos del mercado cambiario.
\end{verbatim}

\section{Paso 12: Crear consulta de monitoreo}

Crear:

\begin{verbatim}
sql/monitor_staging.sql
\end{verbatim}

Contenido:

\begin{lstlisting}[language=SQL]
SELECT
COUNT(*) AS total_filas
FROM staging.tipo_cambio_raw;

SELECT
COUNT(*) AS total_filas
FROM staging.balance_bancario_raw;
\end{lstlisting}

Esta consulta se utilizará para monitorear futuras cargas.

\section{Buenas prácticas}

Durante este proyecto seguiremos las siguientes reglas:

\begin{enumerate}
\item Nunca modificar datos manualmente en staging.
\item Mantener los datos lo más cercanos posible a la fuente.
\item No realizar cálculos en staging.
\item Todas las transformaciones ocurren en la capa core.
\item Conservar metadatos de carga.
\end{enumerate}

\section{Checklist de validación}

Antes de continuar verificar:

\begin{itemize}
\item[$\Box$] Schema staging creado.
\item[$\Box$] Schema core creado.
\item[$\Box$] Schema mart creado.
\item[$\Box$] Tabla tipo_cambio_raw creada.
\item[$\Box$] Tabla balance_bancario_raw creada.
\item[$\Box$] Scripts SQL almacenados en la carpeta sql.
\item[$\Box$] Consultas de validación ejecutadas.
\end{itemize}

\section{Entregable}

Al finalizar este hito el estudiante deberá disponer de:

\begin{itemize}
\item Arquitectura física inicial del Data Warehouse.
\item Schemas creados.
\item Tablas staging creadas.
\item Scripts SQL documentados.
\item Capa Bronze lista para recibir datos.
\end{itemize}

\section{Próximo hito}

En el Hito 6 construiremos el primer proceso ETL.

Aprenderemos a:

\begin{itemize}
\item Leer archivos CSV desde Python.
\item Conectarnos a PostgreSQL.
\item Cargar información a staging.
\item Automatizar la ingesta.
\item Validar que los datos fueron cargados correctamente.
\end{itemize}
